%% kapa: exam=1 round=31 year=2020 period=2 form=A subject=appraisal_law subject_ko=감정평가관계법규 seq=3 number=3
%% source: Q-Net 감정평가사 기출문제 (cst00309) artlSeq=5207219 file=extracted/1008/5207219/A형/A형/2020년도 제31회 감정평가사 1차 시험_2교시_A형.pdf
%% conversion: kordoc 4.24.51 -> tools/gen1.py (auto); review=proofread-sonnet; flags=none
\begin{question}{3}
국토의 계획 및 이용에 관한 법령상 개발행위의 허가 등에 관한 설명으로 옳은 것은?
\begin{choices}
  \item 재난수습을 위한 응급조치인 경우에도 개발행위허가를 받고 하여야 한다.
  \item 시장 또는 군수가 개발행위허가에 경관에 관한 조치를 할 것을 조건으로 붙이는 경우 미리 개발행위허가를 신청한 자의 의견을 들어야 한다.
  \item 성장관리방안을 수립한 지역에서 하는 개발행위는 중앙도시계획위원회와 지방도시계획위원회의 심의를 거쳐야 한다.
  \item 지방자치단체는 자신이 시행하는 개발행위의 이행을 보증하기 위하여 이행보증금을 예치하여야 한다.
  \item 기반시설부담구역으로 지정된 지역은 중앙도시계획위원회의 심의를 거쳐 10년 이내의 기간 동안 개발행위허가를 제한할 수 있다.
\end{choices}
\end{question}
\answer{2}
